\section{Sequence Diagrams}
% ══════════════════════════════════════════════════════════════════════════════
% 0.1
% ══════════════════════════════════════════════════════════════════════════════
\subsection{ Admin Authentication \& Dashboard Sequence}
\begin{center}
	\includegraphics[width=\textwidth]{diagrams/seq-01.pdf}
	\captionof{figure}{Admin Authentication \& Dashboard Sequence}
	\label{fig:admin-auth-seq}
\end{center}

This Sequence Diagram illustrates the interaction between the Administrator, the System, and the Database (DBMS) during the authentication process and the
subsequent loading of the dashboard.\textbf{Key Technical Logic:}

\begin{itemize}
	\item \textbf{Authentication Flow}
	\begin{itemize}
		\item The process begins with the Admin submitting credentials.
		\item The System validates these credentials against the DBMS.
		\item An \textbf{Alternative Frame (alt)} handles two scenarios:
		\begin{itemize}
			\item \textbf{Success:} The user is authenticated and granted access to the system.
			\item \textbf{Failure:} An error message is triggered and access is denied.
		\end{itemize}
	\end{itemize}
	
	\item \textbf{Data Retrieval (Dashboard Request)}
	\begin{itemize}
		\item Upon successful login, the Admin requests the Dashboard.
		\item The System executes a query to the DBMS to \textbf{fetch statistics} (e.g., total patients, appointments, revenue).
	\end{itemize}
	
	\item \textbf{Data Processing and Presentation}
	\begin{itemize}
		\item After receiving the raw data, the System processes it (\textbf{Load Dashboard Statistics}).
		\item The System renders the final visual dashboard interface for the Admin.
	\end{itemize}
\end{itemize}
% ══════════════════════════════════════════════════════════════════════════════
% 0.2
% ══════════════════════════════════════════════════════════════════════════════

\subsection{Admin Authentication with Logging}
\addcontentsline{toc}{subsection}{0.2 Admin Authentication with Logging}

\begin{figure}[H]
	\centering
	\includegraphics[width=0.82\textwidth]{diagrams/seq-02.pdf}
	\caption{Admin Authentication with Logging}
	\label{fig:admin-auth-logging}
\end{figure}

\textbf{Description:} \textit{This Sequence Diagram extends the basic authentication
	flow by integrating a Security Audit Mechanism. It demonstrates how the system tracks
	access attempts for security monitoring and compliance.}

\medskip
\textbf{Key Technical Logic:}

\begin{itemize}
	\item \textbf{Audit Integration:} Unlike the standard login, this flow includes a
	mandatory interaction with the \texttt{Active\_log} (via the DBMS) inside the
	alternative block (\texttt{alt}).
	
	\item \textbf{Security Logging Scenarios:}
	\begin{itemize}
		\item \textbf{Success Logging:} If credentials are valid, the system
		triggers a \texttt{Log(action: LOGIN\_SUCCESS)} before granting
		access, creating a permanent record of the session.
		\item \textbf{Failure Logging:} In case of incorrect credentials, the
		system records a \texttt{Log(action: LOGIN\_FAILED)}. This is a
		critical feature to detect Brute Force attacks.
	\end{itemize}
	
	\item \textbf{Data Integrity:} The logging action is performed directly on the
	Database level, ensuring that even if the user session ends or is interrupted,
	the security record remains intact.
\end{itemize}

\subsection{Admin User Management Sequence}
\addcontentsline{toc}{subsection}{0.3 Admin User Management Sequence}

\begin{figure}[H]
	\centering
	\includegraphics[width=0.82\textwidth]{diagrams/seq-03.pdf}
	\caption{Admin User Management Sequence}
	\label{fig:admin-user-mgmt}
\end{figure}

\textbf{Description:} \textit{This Sequence Diagram details the operational flow for
	managing system users, including Doctors and Staff members. It illustrates the
	administrative control over the system's human resources through two primary
	workflows: Creation and Modification.}

\medskip
\textbf{Key Technical Logic:}

\begin{itemize}
	\item \textbf{Workflow Selection:} The process begins with the Administrator
	selecting a specific user category, which triggers the system to provide the
	relevant management interface via a dynamic Control Panel.
	
	\item \textbf{Administrative Operations (\texttt{alt} frame):}
	\begin{itemize}
		\item \textbf{Creation Flow (Add New User):} This branch demonstrates
		initializing a new record, submitting validated data, and committing
		it to the DBMS. The system provides immediate feedback upon
		successful insertion.
		\item \textbf{Modification Flow (Edit / Deactivate):} This covers the
		lifecycle of existing records. It allows the admin to retrieve current
		details and modify fields or change the user's status to
		\textit{Deactivated} (Soft Delete) to maintain historical data
		integrity.
	\end{itemize}
	
	\item \textbf{Data Persistence:} All operations culminate in a confirmation from
	the DBMS, ensuring that the user interface is always synchronized with the
	actual state of the database, preventing Data Inconsistency.
\end{itemize}

\newpage

% ══════════════════════════════════════════════════════════════════════════════
% 0.4
% ══════════════════════════════════════════════════════════════════════════════
\subsection{ Admin User Management with Logging}
\addcontentsline{toc}{subsection}{0.4 Admin User Management with Logging}

\begin{figure}[H]
	\centering
	\includegraphics[width=0.82\textwidth]{diagrams/seq-02.pdf}
	\caption{Admin User Management with Logging}
	\label{fig:admin-user-mgmt-logging}
\end{figure}

\textbf{Description:} \textit{This Sequence Diagram illustrates the Secure Resource
	Provisioning flow. It emphasizes the mandatory audit step required for sensitive
	administrative actions, such as creating new user accounts, ensuring every transaction
	is fully traceable.}

\medskip
\textbf{Key Technical Logic:}

\begin{itemize}
	\item \textbf{Transactional Integrity:} The process follows a strict sequence where
	user data is first persisted in the primary user table. Only after receiving a
	\textbf{Save Confirmation} from the DBMS does the system proceed to the
	next phase.
	
	\item \textbf{Traceability (Security Logging):} A secondary interaction with the
	DBMS is performed to execute the \texttt{Log(action: CREATE\_USER)}
	command. This ensures that every new account is linked to the specific
	Administrator, providing a clear Audit Trail for security monitoring.
	
	\item \textbf{Sequential Synchronization:} To prevent ``unlogged actions,'' the
	Success Message is only delivered to the Admin after both the primary data
	saving and the secondary logging are successfully confirmed by the database.
\end{itemize}

% ══════════════════════════════════════════════════════════════════════════════
% 0.5
% ══════════════════════════════════════════════════════════════════════════════
\subsection*{0.5 \quad Admin Report Generation and Export}
\addcontentsline{toc}{subsection}{0.5 Admin Report Generation and Export}

\begin{figure}[H]
	\centering
	\includegraphics[width=0.82\textwidth]{diagrams/seq-05.pdf}
	\label{fig:admin-report-export}
\end{figure}

\textbf{Description:} \textit{This Sequence Diagram illustrates the lifecycle of report
	generation and the secure export of clinical or financial data. It highlights the system's
	ability to transform raw database records into portable document formats while
	maintaining a strict audit trail for data privacy.}

\medskip
\textbf{Key Technical Logic:}

\begin{itemize}
	\item \textbf{Data Aggregation:} The process begins with the Administrator's
	request, triggering the System to \textit{``Fetch Data''} from the DBMS.
	This step involves complex queries to aggregate statistics or detailed patient
	history.
	
	\item \textbf{Internal Processing (PDF Generation):} A self-invocation message
	\textit{Generate PDF} represents the logic layer where the system (utilizing
	PHP libraries) formats the raw data into a structured, read-only document.
	
	\item \textbf{Security \& Compliance:} Before the file is provided, the system
	executes a \texttt{Log(action: REPORT\_EXPORT)}. This ensures that any
	data leaving the system is tracked, satisfying critical requirements for medical
	data privacy and HIPAA-style standards.
	
	\item \textbf{Secure Delivery:} The final \textit{``Export File''} is delivered to
	the Admin only after the logging action is confirmed. This ensures full
	accountability for the distribution of sensitive clinical or financial information.
\end{itemize}



% ══════════════════════════════════════════════════════════════════════════════
% 0.6
% ══════════════════════════════════════════════════════════════════════════════
\subsection{Admin Settings Update with Logging}
\addcontentsline{toc}{subsection}{0.6 Admin Settings Update with Logging}

\begin{figure}[H]
	\centering
	\includegraphics[width=0.82\textwidth]{diagrams/seq-06.pdf}
	\caption{Admin Settings Update with Logging}
	\label{fig:admin-settings-logging}
\end{figure}

\textbf{Description:} \textit{This Sequence Diagram details the flow for modifying
	global system configurations. Given the high impact of these settings on the
	application's behavior, the process integrates an immediate auditing step to ensure
	transparency and stability.}

\medskip
\textbf{Key Technical Logic:}

\begin{itemize}
	\item \textbf{Configuration Persistence:} When the Admin submits updated
	settings, the System interacts directly with the \texttt{Config/Settings}
	table. This ensures that global parameters (such as clinic hours, contact
	info, or system thresholds) are applied instantly across the entire platform.
	
	\item \textbf{Audit Enforcement:} Following a successful update, the System
	triggers a \texttt{Log(action: UPDATE\_SETTINGS)} entry. This record serves
	as a critical reference point for troubleshooting or security reviews if system
	behavior is altered.
	
	\item \textbf{Final Confirmation:} The Administrator receives a
	\textit{``Settings Saved''} notification only after the database confirms both
	the configuration update and the security log entry, ensuring total
	synchronization between the UI and the actual System State.
\end{itemize}


% ══════════════════════════════════════════════════════════════════════════════
% 0.7
% ══════════════════════════════════════════════════════════════════════════════
\subsection{Reports and System Settings Sequence}
\addcontentsline{toc}{subsection}{0.7 Reports and System Settings Sequence}

\begin{figure}[H]
	\centering
	\includegraphics[width=0.82\textwidth]{diagrams/seq-07.pdf}
	\caption{Admin Reports \& Settings Sequence}
	\label{fig:admin-reports-settings}
\end{figure}

\textbf{Description:} \textit{This Sequence Diagram provides a consolidated view of
	two critical administrative functions: Data Analytics and System Configuration. It
	demonstrates how the system handles diverse requests through conditional logic to
	maintain operational efficiency.}

\medskip
\textbf{Key Technical Logic:}

\begin{itemize}
	\item \textbf{Reporting Workflow (Data Analytics):}
	\begin{itemize}
		\item \textbf{Parameter Definition:} The Administrator initiates the
		process by defining specific parameters such as Report Type and
		Date Range.
		\item \textbf{Visual Transformation:} The System fetches the relevant
		dataset from the DBMS and performs an internal process
		(\textit{Generate Graphical Report}) to transform raw data into
		visual formats like charts or tables.
		\item \textbf{Optional Export:} An optional \textit{``Export''} step
		allows the user to download the final document for offline use or
		archiving.
	\end{itemize}
	
	\item \textbf{Settings Management Workflow:}
	\begin{itemize}
		\item \textbf{Dynamic Updates:} This section covers the update of clinic
		metadata and operational constraints through two distinct
		interactions: updating general information and modifying
		organizational rules (Working Hours).
		\item \textbf{Real-time Synchronization:} Each action triggers a direct
		update to the DBMS, ensuring that clinical availability and
		system-wide information are always current and synchronized across
		the platform.
	\end{itemize}
\end{itemize}


\newpage

% ══════════════════════════════════════════════════════════════════════════════
% 0.8
% ══════════════════════════════════════════════════════════════════════════════
\subsection{Doctor Consultation and Medical Record Creation}
\addcontentsline{toc}{subsection}{0.8 Doctor Consultation and Medical Record Creation}

\begin{figure}[H]
	\centering
	\includegraphics[width=\textwidth]{diagrams/seq-08.pdf}
	\caption{Doctor Consultation and Medical Record Creation}
	\label{fig:doctor-consultation}
\end{figure}

\textbf{Description:} \textit{This Sequence Diagram illustrates the critical clinical
	workflow, divided into two primary stages: information retrieval and data persistence.
	It showcases the seamless interaction between the Doctor, the Vue.js Interface, the
	PHP Controller, and the Database to ensure high-quality patient care.}

\medskip
\textbf{Key Technical Logic:}

\begin{itemize}
	\item \textbf{Stage 1 -- Historical Review (Data Retrieval):} The process begins
	with the Doctor retrieving the patient's medical history. The System executes
	a secure query to the DBMS to fetch all previous records, ensuring the Doctor
	has the necessary clinical context before the examination begins.
	
	\item \textbf{Stage 2 -- Secure Data Entry \& Validation:} Upon entering the
	diagnosis and treatment plans, the \texttt{Consultation Controller} performs
	an internal validation check to ensure data consistency before attempting
	persistence.
	
	\item \textbf{Transaction Management (Data Integrity):} To ensure total
	integrity, the system utilizes a Database Transaction
	(\texttt{BEGIN / COMMIT}). The medical record is inserted, and a security
	log is simultaneously generated as part of a single unit of work.
	
	\item \textbf{Atomic Operations:} If any part of the process fails, a
	\texttt{ROLLBACK} is triggered to prevent partial or corrupted data. A
	Success Message is only displayed after a successful commit, guaranteeing
	that the medical history is accurately and permanently updated.
\end{itemize}



\newpage

% ══════════════════════════════════════════════════════════════════════════════
% 0.9
% ══════════════════════════════════════════════════════════════════════════════
\subsection{Guest Appointment Booking (No Login)}
\addcontentsline{toc}{subsection}{0.9 Guest Appointment Booking (No Login)}

\begin{figure}[H]
	\centering
	\includegraphics[width=\textwidth]{diagrams/seq-09.pdf}
	\caption{Guest Appointment Booking (No Login)}
	\label{fig:guest-booking}
\end{figure}

\textbf{Description:} \textit{This Sequence Diagram details the workflow for
	spontaneous booking by guest patients. It emphasizes real-time availability checking
	and atomic data handling to ensure a seamless and secure entry point into the system
	for unauthenticated users.}

\medskip
\textbf{Key Technical Logic:}

\begin{itemize}
	\item \textbf{Availability Management (Real-time):} The process starts with a
	dynamic fetch of available slots from the DBMS. The system filters records in
	real-time to ensure only \textbf{``Free''} slots are displayed, preventing
	overbooking directly at the UI level.
	
	\item \textbf{Guest Data Validation:} Since the user is not authenticated, the
	\texttt{Booking Controller} performs rigorous server-side validation on the
	provided contact information (Name, Phone, National ID) to maintain high
	data quality and integrity.
	
	\item \textbf{Resource Reservation (Race Condition Prevention):} To prevent
	Race Conditions -- where multiple guests might attempt to book the same slot
	simultaneously -- the system wraps the appointment insertion and the slot
	status update within a strict Database Transaction.
	
	\item \textbf{Output Generation \& Bridging:} Upon a successful \texttt{COMMIT},
	the system generates a unique Booking Slip (PDF/QR Code). This serves as a
	temporary credential, effectively bridging the gap between an anonymous
	guest and a formally registered patient record upon arrival.
\end{itemize}
\newpage

% ══════════════════════════════════════════════════════════════════════════════
% 0.10
% ══════════════════════════════════════════════════════════════════════════════
\subsection*{0.10 \quad View Medical Records and Results (Post-Registration)}
\addcontentsline{toc}{subsection}{0.10 View Medical Records and Results (Post-Registration)}

\begin{figure}[H]
	\centering
	\includegraphics[width=\textwidth]{diagrams/seq-10.pdf}
	\caption{View Medical Records and Results (Post-Registration)}
	\label{fig:patient-view-records}
\end{figure}

\textbf{Description:} \textit{This Sequence Diagram illustrates the secure process for
	registered patients to access their clinical data via the Patient Portal. It emphasizes a
	multi-layered approach to authentication and data-level security to protect sensitive
	medical history.}

\medskip
\textbf{Key Technical Logic:}

\begin{itemize}
	\item \textbf{Secure Access (Authentication):} The workflow begins with a
	mandatory authentication phase. Patients utilize secure credentials provided
	by the clinic staff, ensuring that only verified individuals can access the
	portal's services.
	
	\item \textbf{Access Control (Session Filtering):} A critical security measure is
	implemented where the \texttt{Records Controller} filters the database query
	to ensure the \texttt{patient\_id} strictly matches the active session. This
	mechanism prevents Horizontal Privilege Escalation and unauthorized access
	to other patients' records.
	
	\item \textbf{Audit Logging (Confidentiality):} Each time a patient views their
	records, a \texttt{Log(action: VIEW\_RECORDS)} is generated. This maintains
	a complete history of data access, ensuring compliance with medical
	confidentiality standards and accountability.
	
	\item \textbf{Information Delivery:} The system provides a structured list of
	medical histories and lab results. The interface allows patients to view or
	download detailed reports in PDF format for personal archiving or sharing
	with other healthcare providers.
\end{itemize}

\newpage

% ══════════════════════════════════════════════════════════════════════════════
% 0.11
% ══════════════════════════════════════════════════════════════════════════════
\subsection{Online Payment and Invoicing}
\addcontentsline{toc}{subsection}{0.11 Online Payment and Invoicing}

\begin{figure}[H]
	\centering
	\includegraphics[width=\textwidth]{diagrams/seq-11.pdf}
	\caption{Online Payment and Invoicing}
	\label{fig:patient-payment}
\end{figure}

\textbf{Description:} \textit{This Sequence Diagram details the financial transaction
	workflow within the Patient Portal. It illustrates the seamless integration between the
	clinic's internal logic and external Payment Gateways, ensuring secure and consistent
	billing operations.}

\medskip
\textbf{Key Technical Logic:}

\begin{itemize}
	\item \textbf{Invoicing Retrieval (Targeted Billing):} The process begins with the
	System fetching unpaid invoices from the DBMS. This ensures that the Patient
	only interacts with relevant billing data strictly associated with their specific
	appointments and records.
	
	\item \textbf{External Gateway Integration (API Handling):} When the Patient
	initiates payment, the \texttt{Payment Controller} communicates with an
	external Payment Gateway (Bank) using secure tokens. This highlights the
	system's architecture for handling external API responses and managing
	asynchronous transaction statuses.
	
	\item \textbf{Atomic Financial Updates (Multi-table Consistency):} Upon a
	successful transaction, the system utilizes a Database Transaction to
	perform three simultaneous actions as a single unit:
	\begin{itemize}
		\item Updating the payment status to \texttt{'Paid'}.
		\item Closing the billing cycle for the associated appointment.
		\item Recording a security log (\texttt{ONLINE\_PAYMENT}) for financial
		auditing.
	\end{itemize}
	
	\item \textbf{Success \& Feedback (Legal Proof):} The workflow concludes with
	the generation of a digital Receipt (PDF), providing the patient with
	immediate legal proof of payment. The use of \texttt{COMMIT / ROLLBACK}
	logic guarantees that no partial financial updates occur, maintaining absolute
	data consistency.
\end{itemize}

\newpage

% ══════════════════════════════════════════════════════════════════════════════
% 0.12
% ══════════════════════════════════════════════════════════════════════════════
\subsection{Patient Registration Sequence}
\addcontentsline{toc}{subsection}{0.12 Patient Registration Sequence}

\begin{figure}[H]
	\centering
	\includegraphics[width=\textwidth]{diagrams/seq-12.pdf}
	\caption{Patient Registration Sequence (Receptionist Module)}
	\label{fig:patient-registration}
\end{figure}

\textbf{Description:} \textit{This Sequence Diagram illustrates the standardized
	procedure for onboarding new patients into the clinic's database. It emphasizes data
	integrity and auditability through a rigorous two-step verification and creation process.}

\medskip
\textbf{Key Technical Logic:}

\begin{itemize}
	\item \textbf{Duplicate Prevention (Integrity Check):} The workflow begins with
	a mandatory search for existing records using unique identifiers (National ID
	or Phone). This prevents the creation of redundant profiles, ensuring that
	each patient maintains a single, consistent \textbf{Master Medical Record}.
	
	\item \textbf{Conditional Workflow (\texttt{alt} frame):} The system utilizes an
	Alternative Frame (\texttt{alt}) to handle dynamic scenarios:
	\begin{itemize}
		\item \textbf{Existing Patient Found:} The system bypasses registration
		and retrieves the current profile, preventing data fragmentation.
		\item \textbf{New Patient Registration:} The system enables the entry
		form only after confirming the absence of a duplicate record in
		the DBMS.
	\end{itemize}
	
	\item \textbf{Secure Persistence \& Accountability:} During registration, the
	\texttt{Patient Controller} manages the insertion of personal details.
	Simultaneously, a critical security log
	\texttt{Log(action: REGISTER\_PATIENT)} is recorded. This links the new
	record to the specific receptionist, ensuring full \textbf{Traceability} for
	all administrative actions.
\end{itemize}

\newpage
% ══════════════════════════════════════════════════════════════════════════════
% 0.13
% ══════════════════════════════════════════════════════════════════════════════
\subsection{Appointment Booking And Availability}
\addcontentsline{toc}{subsection}{0.13 Appointment Booking And Availability}

\begin{figure}[H]
	\centering
	\includegraphics[width=\textwidth]{diagrams/seq-14.pdf}
	\caption{Appointment Booking And Availability}
	\label{fig:appointment-booking}
\end{figure}

\textbf{Description:} \textit{This Sequence Diagram illustrates a robust workflow for
	manual appointment scheduling by the Receptionist. It is designed to ensure maximum
	data consistency and prevent double-booking through rigorous real-time validation.}

\medskip
\textbf{Key Technical Logic:}

\begin{itemize}
	\item \textbf{Real-time Availability (Schedule View):} The system queries the
	DBMS to identify existing appointments for a specific doctor and date. The
	\texttt{Appointment Controller} then dynamically calculates free slots,
	providing the Receptionist with an accurate, real-time visual representation
	of the schedule.
	
	\item \textbf{Conflict Prevention (Validation Step):} Upon selecting a slot, the
	system initiates a Database Transaction. It performs a critical secondary
	\texttt{SELECT status} check to ensure the slot hasn't been reserved by
	another user (e.g., via the Patient Portal) in the interval between viewing
	and confirming.
	
	\item \textbf{Atomic Commitment \& Concurrency:}
	\begin{itemize}
		\item \textbf{Success:} If the slot remains available, the record is
		inserted, the transaction is \texttt{COMMITTED}, and a security
		log is generated.
		\item \textbf{Failure (Rollback):} If a conflict is detected, the
		transaction is \texttt{ROLLED BACK} immediately, protecting the
		database from corrupted or overlapping schedules.
	\end{itemize}
	
	\item \textbf{User Feedback:} The process concludes with either a printable
	Appointment Slip for the patient or a mandatory schedule refresh in case of
	failure, ensuring the UI is always synchronized with the DBMS state.
\end{itemize}

\newpage

% ══════════════════════════════════════════════════════════════════════════════
% 0.14
% ══════════════════════════════════════════════════════════════════════════════
\subsection{Billing And Receipt Generation}
\addcontentsline{toc}{subsection}{0.14 Billing And Receipt Generation}

\begin{figure}[H]
	\centering
	\includegraphics[width=\textwidth]{diagrams/seq-14.pdf}
	\caption{Billing And Receipt Generation}
	\label{fig:billing-receipt}
\end{figure}

\textbf{Description:} \textit{This Sequence Diagram illustrates the final stage of the
	patient visit: the transition from clinical service to financial settlement. It details how
	the system ensures accurate billing and secure payment recording through automated
	calculations and transactional integrity.}

\medskip
\textbf{Key Technical Logic:}

\begin{itemize}
	\item \textbf{Dynamic Invoice Calculation:} The process begins by fetching
	service details and fees associated with the completed appointment. The
	\texttt{Billing Controller} performs an internal calculation
	(\texttt{Total = Services + Taxes}), ensuring the Receptionist presents an
	accurate \textbf{``Invoice Draft''} before any transaction occurs.
	
	\item \textbf{Secure Payment Processing (Integrity):} Upon receiving payment,
	the system initiates a Database Transaction. It records the payment details
	and simultaneously updates the \texttt{billing\_status} to
	\texttt{'Closed'}. This dual action ensures no appointment remains
	``Unpaid'' in the records after a successful transaction.
	
	\item \textbf{Accountability \& Audit Trail:} A security log
	\texttt{Log(action: GENERATE\_BILL)} is created to track the financial
	transaction. This is critical for Daily Reconciliation and internal audits,
	linking the revenue to the specific receptionist and appointment.
	
	\item \textbf{Document Output \& Finalization:} The workflow concludes with the
	generation of a Digital Receipt (PDF). This provides a professional,
	printable document for the patient while maintaining a permanent, immutable
	digital copy within the system for future reference.
\end{itemize}
